\chapter{绪论}
\section{模板用途与适用范围}
本模板依据长沙理工大学2023年研究生学位论文撰写规范整理，首版适用于全日制专业硕士。学术硕士、博士、非全日制及本科论文的封面和学位信息有所不同，不应仅修改论文类别文字后直接套用。

模板的目标是减少重复的格式操作。正文默认采用小四宋体，英文字母与数字采用Times New Roman，行距为1.5倍；各级标题、图表题注和参考文献分别拥有独立的格式设置。论文结构和研究内容仍应由作者根据培养要求自行完成。
\section{文件组织与写作流程}
个人信息统一填写在配置文件中。论文正文按章放入章节目录，文献信息保存在BibTeX数据库，文献引用由BibTeX与natbib处理。参考文献的顺序编码格式使用GB/T 7714样式\cite{gbt7714}。
\subsection{开始写作}
先编译未经修改的示例，确认环境和字体可用；然后替换题目、作者和中英文摘要。新增章节时在主文件中加入对应文件，避免复制整套排版配置。
\subsubsection{避免局部覆盖}
不要为凑页数而随意使用字号缩小命令或改变正文行距。长公式应按语义换行，表格应调整列宽，图形应保留足够清晰的文字。使用书籍介绍的基本LaTeX命令，可以保持源文件结构清楚\cite{lamport1994}。
\section{章节安排}
第二章演示公式、表格与插图的用法，并展示跨页表格和子图题注。第三章说明如何替换示例、完成提交前检查及反馈格式问题。示例的目的在于展示可复用结构，并不构成真实论文内容。
